\begin{abstract}
Low-cost potentiostats are routinely validated against a resistive dummy cell by
reporting the coefficient of determination of an Ohm's-law fit. We show that
this evidence is structurally incapable of detecting a transimpedance gain
error: a constant gain multiplies the fitted slope and leaves $R^{2}$ exactly
unchanged. The point is not hypothetical. An AD5941/RP2040 potentiostat examined
here recovered a resistance $424\,\%$ in error from cyclic voltammetry---its
reported currents low by a factor of $5.22$---while fitting
Ohm's law at $R^{2} = 0.99999647$ against a precision Randles cell
($R_s = \SI{560}{\ohm}$, $R_{ct} = \SI{10}{\kilo\ohm}$, $C_{dl} = \SI{33}{\nano\farad}$).
The cause was a single firmware constant declaring the value of an on-board
calibration resistor as \SI{10}{\kilo\ohm} when the resistor fitted was
\SI{200}{\ohm}; because the front end measures its own transimpedance gain
against that declaration, the error entered every reported current
multiplicatively. Deliberately sweeping the declared value moved the absolute
error by a factor of $982$ while $R^{2}$ varied by $2.4\times10^{-6}$. The error
also proved non-correctable after the fact, because the calibration excitation is
itself sized from the declared value and saturates. We propose a five-check
validation protocol---scale traceability, absolute accuracy against component
nominals, residual structure referred to the current quantisation step,
replicate precision, and invariance across acquisition settings---that requires
no reference instrument. Applied after correcting eight firmware defects, it
placed the instrument's cyclic voltammetry at $+0.428\,\%$ of the nominal
\SI{10560}{\ohm} reference with a replicate relative standard deviation of
$0.0152\,\%$, invariant to within $0.074$ percentage points across scan rates of
\SIrange{25}{800}{\milli\volt\per\second}, potential limits to
$\pm\SI{600}{\milli\volt}$ and steps of \SIrange{5}{20}{\milli\volt}, with fit
residuals at $0.25$ of the quantisation step. We also report what the protocol
shows about the instrument's five other implemented methods: chronoamperometry,
square-wave and differential pulse voltammetry return no cell current,
impedance spectroscopy completes no frequency point, and open-circuit potential
saturates. One method of six is validated. The principal limitations are that
traceability rests on component nominal values rather than on a calibrated
standard, that a single passive operating point at a single transimpedance
setting was exercised, and that the cell's capacitance lies below the
instrument's current resolution and so was not recovered.
\end{abstract}

\begin{keyword}
potentiostat \sep instrument validation \sep calibration traceability \sep
cyclic voltammetry \sep AD5941 \sep dummy cell \sep open-source hardware
\end{keyword}
